\section{用户洞察与问题定义}

\subsection{目标用户:三档 KOC 画像}

我们调研了行业研报(克劳锐《2024 年 KOL 发展年报》)、深度报道(36 氪「郑州帮」)、知乎/CSDN/B 站
真实创作者自述,把目标用户精确划分为三档,Ripple 为其中两档(B/C)提供差异化价值。

\subsubsection{画像 A:兼职素人 KOC(1k--5k 粉,月收入 0--5000 元)}

\begin{quote-box}
\textbf{典型一天}\\
早晨地铁通勤刷小红书 / 抖音找「能抄结构」的爆款链接;白天主业间隙写口播大纲;
晚上 1.5--3 小时剪映粗剪 + 自动字幕 + 模板封面;睡前看后台数据,
\textbf{常常归因于「限流」或「今天运气不好」}。
\end{quote-box}

\textbf{核心特征}:
\begin{itemize}
\item 工具预算 ¥0--200/月,免费党为主
\item 工具支出 95\% 在剪映
\item 「自然量过线再投 DOU+」是民间纪律
\item 接受「内容方向建议」但不愿付订阅
\end{itemize}

\textbf{Ripple 对画像 A 的价值}:免费层的「灵感库 + 早期信号订阅(每天 3 条)」,
覆盖兼职素人,是用户量的基本盘。

\subsubsection{画像 B:腰部 KOC / 小团队(5k--5万粉,月收入 5000--2万)}

\begin{quote-box}
\textbf{典型一天}\\
上午看昨日笔记 / 视频赞藏与「\textbf{搜索占比}」(尤其小红书);
中午追热点快反;下午拍摄外出探店;晚间定稿标题、封面 A/B 思路,
考虑是否投薯条 / DOU+;商业稿件改稿轮次与排期撞车是常见痛点。
\end{quote-box}

\textbf{核心特征}:
\begin{itemize}
\item 工具预算 ¥100--800/月,愿意为「省时间」付费
\item 已有微弱团队(摄影 / 剪辑 / 助理)
\item 商业合作开始变多,蒲公英 / 星图报价频繁
\item 持续焦虑「哪个选题能爆」
\end{itemize}

\textbf{Ripple 对画像 B 的价值}:Pro 版的「12 Agent 工厂 + 早期信号雷达 + 风格记忆」,
ARPU ¥99/月,是商业化的主力。

\subsubsection{画像 C:机构化「素人矩阵」执行者(MCN / 品牌侧 KOC 团队)}

\begin{quote-box}
\textbf{36 氪原文引用}\\
「\textbf{竞品背后有 80 人规模的小红书团队在支撑},才终于恍然大悟:
『难怪打不过。』」

「按模板拍照、修图、写文案,\textbf{下班前发满规定条数(如 6 条)};
发布后盯数据,破百赞要马上补评、引导到站外购买链接。」
\end{quote-box}

\textbf{核心特征}:
\begin{itemize}
\item 工具预算 ¥数千 / 月,B 端采购流程
\item 多账号统一管理,需协同与权限
\item 关心 ROI 而非创意本身
\item 决策链长(主管 + IT + 法务)
\end{itemize}

\textbf{Ripple 对画像 C 的价值}:Enterprise 版「多账号管理 + 团队协同 + SLA + 私有部署」,
ARPU 5--20 万 / 年,是高 ARPU 故事的支撑。

% --------------------------------------------

\subsection{用户痛点:十大不可言说的真实痛点}

我们刻意避开「内容质量不稳」「定位不清」这种官方话术,
深入 \textbf{36 氪深度报道、知乎收益讨论、CSDN 技术博客、B 站 Vlog 自述},
提取 KOC 真实但不愿明说的痛点。每条标注来源,可追溯。

\begin{enumerate}[leftmargin=*]

\item \textbf{「认真写长文不如抖机灵」}\\
\textit{原话(网易转载知乎)}:「认真运营的创作者,一篇文章写一两天,产出篇数少,
以同样的权重,去和一分钟写一条的回答争抢流量,难度太大。」\\
\textit{Ripple 解法}:Oracle Agent 的 \textbf{早期信号雷达} 让长文创作者
也能精准押中 7--14 天后的爆点,避免与短平快内容拼数量。

\item \textbf{「10W 阅读分到 26 元」收益绝望}\\
\textit{原话(同上)}:「一篇 10W 阅读的文章,实际收益是多少?
\textbf{可怜的 26.18 元},只有公众号收益的 455 分之一。」\\
\textit{Ripple 解法}:多平台适配 + 跨平台分发,一稿多吃,
帮 KOC 把同一选题在公众号(高 RPM)+ 视频号(分成)+ 小红书(种草接单)同时变现。

\item \textbf{「播放量卡 2 万像追女神」算法玄学焦虑}\\
\textit{原话(运营喵)}:「播放量卡在 2 万就像追女神\dots」\\
\textit{Ripple 解法}:InsightAnalystAgent 提供 \textbf{归因解释}——
不是给玄学归因,而是基于早期信号 + 仿真数据告诉用户「这条预测应该会爆 X 概率,
若没爆是因为 Y(发布时间 / 钩子 / 封面)。」

\item \textbf{「MCN 只有接广告时才是有效沟通」}\\
\textit{原话(36 氪)}:「\textbf{当自己没有选题灵感时,对方给出的方案十分敷衍},
唯一的有效沟通,就是接广告,且拒接会有违约风险。」\\
\textit{Ripple 解法}:把 MCN 的「灵感支持职能」用 AI Agent 替代,
KOC 不再依赖人际关系拿到选题。

\item \textbf{「单条广告 8 万只分 500」分账不公}\\
\textit{原话(行业媒体)}:网友爆料「一单价格 4.8 万--8.5 万的商单,
\textbf{隋坡仅分成 500 元}」引发讨论。\\
\textit{Ripple 解法}:虽然不直接解决分账问题,但帮 KOC 提升「自主接单能力」,
脱离 MCN 也能持续产出。

\item \textbf{「降 AI 工具把『绝』改成『卓越』被评论区喷装」}\\
\textit{原话(CSDN/GitCode)}:「嘎嘎降 AI 把我一句『这粉底液持妆真的绝』改成了
『此粉底之持妆效能确系卓越』\dots\textbf{评论区直接有人喷我装}。」\\
\textit{Ripple 解法}:\texttt{ai-detection-bypass} Skill 内置 15 个具体技巧
(强制写入线下细节、句式长度抖动、口头禅插入位置等),保证「降 AI 味但保留人味」。

\item \textbf{「批量抄爆款像流水线但洗稿不够狠会判搬运」}\\
\textit{原话(高羽网创等)}:「从扒爆款、憋文案、P 封面到发布,手速快也得 15 分钟\dots
\textbf{正文改动低于 30\% 容易被判搬运}」\\
\textit{Ripple 解法}:多 Agent 内容工厂从根上避免「洗稿」——
基于早期信号的原创选题 + 风格学习生成的人化文案,远离搬运灰区。

\item \textbf{「打不过 80 人小红书团队」资源悬殊焦虑}\\
\textit{原话(36 氪)}:「竞品背后有 80 人规模的小红书团队在支撑\dots
\textbf{难怪打不过}。」\\
\textit{Ripple 解法}:12 Agent 协作让单兵 KOC 拥有「8 人小队」的产能(选题 / 写作 /
设计 / 审查 / 复盘 8 个角色都有 AI),抹平资源差距。

\item \textbf{「跨平台一鱼多吃但每个平台风格不同改死人」}\\
\textit{行业共识}:抖音偏「黄金 3 秒强冲击」,小红书偏「首图 + 8 字标题」,
公众号偏「深度长文」,视频号偏「社交链路」——\textbf{同一选题在不同平台需要完全不同的形态}。\\
\textit{Ripple 解法}:ScriptWriterAgent 调 MiniMax M2.7 长上下文,
\textbf{一次性生成 6 个平台的差异化版本}。

\item \textbf{「灵感乍现但来不及记下来,第二天就忘了」}\\
\textit{创作者社群常见抱怨}\\
\textit{Ripple 解法}:\textbf{灵感捕捉系统} ——
macOS Menu Bar / 微信小程序 / iOS Share Extension / Android Quick Settings 全端入口,
\texttt{Cmd+Shift+I} 全局快捷键 → 浮动单行输入 → 异步 AI 处理(转写 / OCR / 自动归档)。

\end{enumerate}

% --------------------------------------------

\subsection{使用场景}

\subsubsection{场景 1:周日晚上 22:30,小红书美妆博主小雯}

\textbf{背景}:小雯,粉丝 1.2 万,做美妆测评 8 个月,月入 8000 元。
明天周一要发新内容但毫无灵感,过去 3 天发的笔记数据都很差(800/1200/600 赞)。

\textbf{原始流程(没有 Ripple)}:
\begin{enumerate}
\item 22:30 打开小红书刷 1 小时找灵感,看到的都是已经爆的选题(红海)
\item 23:30 刷千瓜数据看排行榜,发现「7 天前」的选题已经被同行写了 200 篇
\item 00:30 焦虑入睡,第二天发了一篇模仿爆款的笔记,数据 500 赞
\end{enumerate}

\textbf{Ripple 流程}:
\begin{enumerate}
\item 22:30 打开 Ripple App,点「今日早期信号」
\item 22:32 看到 Top 3 信号:
    \begin{itemize}
    \item 「珂润润浸保湿洗颜泡沫」— 微博讨论量 +180\% / 小红书搜索量 +95\% /
        Polymarket 美妆相关合约升温(置信度 78\%,预计 5--7 天后登榜)
    \item 「冬季敏感肌护理」— 巨量算数趋势 +120\% / 百度指数 +80\%(置信度 65\%)
    \item 「平价彩妆 100 元 5 件」— GitHub Trending 海外类似产品 + 国内代购讨论(置信度 52\%)
    \end{itemize}
\item 22:35 选择 Top 1,Ripple 12 Agent 协作生成完整内容包:
    \begin{itemize}
    \item 10 个候选标题(标注公式:数字 / 否定 / 反差 / 时效)
    \item 5 个候选封面(高对比 / 温和 / 极简 / 人脸 / 无脸)
    \item 笔记正文(1 主图 + 8 副图 + 精炼正文,风格匹配小雯历史作品)
    \item 合规自检通过(无虚假宣传 / 无医疗承诺 / AI 标识就位)
    \item 发布策略建议(明天上午 10:30 发布,粉丝活跃高峰)
\end{itemize}
\item 22:55 小雯审核通过,直接导出发布
\item 第二天 18:00 数据回流:1.8 万赞,涨粉 320,1 个品牌商务联系
\end{enumerate}

\subsubsection{场景 2:周二早上 09:00,公司品牌部经理李总}

\textbf{背景}:李总管理 3 个品牌账号 + 签约 12 个 KOC,月预算 50 万投放。
本周 Q3 季度营销 brief 来了——为「冬季新品保暖羽绒服」找到合适的 KOC + 内容方向。

\textbf{原始流程(没有 Ripple)}:
\begin{enumerate}
\item 9:00 打开飞瓜数据,翻 100+ 美妆生活类 KOC 资料,2 小时只筛出 5 个候选
\item 11:00 给候选 KOC 发 brief 邮件,等回复要 1--3 天
\item 一周后才能定下来内容方向,错过早冬种草窗口
\end{enumerate}

\textbf{Ripple Enterprise 流程}:
\begin{enumerate}
\item 9:00 打开 Ripple Enterprise,输入「冬季羽绒服种草」
\item 9:02 Oracle Agent 输出早期信号 Top 5:
    \begin{itemize}
    \item 「轻量羽绒」+89\% / 「南方湿冷」+62\% / 「通勤穿搭」+48\% \dots
    \end{itemize}
\item 9:05 系统自动从签约 12 个 KOC 中匹配「最适合做这个选题的 3 人」(基于风格库 + 历史数据)
\item 9:10 一键发送「AI 生成的初稿 + 早期信号报告」到 3 个 KOC
\item 当天下午 KOC 回复确认,3 天内全部上线
\item 一周后 Q3 战役提前 4 天启动,GMV 提升 35\%
\end{enumerate}
